\section{Related Work}

\noindent\textbf{Urban perception from street-view imagery.}
Crowdsourced pairwise judgement datasets have enabled vision models
for perceived safety, wealth, beauty, and liveliness at scale~\cite{
salesses2013,naik2014,dubey2016,quintana2025}. These models are
effective for prediction and geographic mapping, but are not designed
to answer mechanistic questions at the scene level: \emph{what visual
evidence drives a judgement, and what localised change would shift it?}

\noindent\textbf{Interpretability and counterfactual explanations.}
Gradient-based saliency and SHAP-style methods attribute predictions
to pixels or semantic features, but are known to be visually plausible
yet unfaithful, failing basic sanity checks~\cite{adebayo2018}.
Concept-based methods such as TCAV~\cite{kim2018} remain correlational
unless paired with explicit interventions. Counterfactual visual
explanations~\cite{goyal2019} and principled causal evaluation
frameworks~\cite{pawlowski2020,melistas2024} provide stronger tools,
but typically target model logits rather than human perceptual
preference, and rarely validate that manipulating identified regions
shifts human judgement under controlled conditions.

\noindent\textbf{Generative editing and urban counterfactuals.}
Diffusion-based editors~\cite{brooks2023,meng2022,zhang2023} enable
photorealistic localised edits, but even constrained edits can
introduce correlated non-target changes (lighting, style, scene
activity) that confound perception evidence. Law et al.~\cite{
law2023,law2025} demonstrate plausible counterfactuals for urban image
regressors and facade design; Zhao et al.~\cite{zhao2026} use
perception-guided generation for streetscape improvement. These
efforts optimise outcomes rather than explain individual scenes, and
do not treat generator artefacts as a threat to validity.

\noindent\textbf{Positioning.}
We unify these threads into a framework for \textbf{interventional
attribution} of street-view perception: explanations are
scene-specific hypotheses tested via constrained prompt-only edits,
audited for non-target drift, and grounded in human pairwise
judgements rather than model logits.
